\section{Solucionadores(solvers)}

En esta seccion veremos 4 solvers para flujo laminar, turbulento, les que esta definido para un numero grande de eddy y otro de dos estados 

\subsection{icoFoam}
Este solver es para fluidos incompresible y laminar, las ecuaciones son
\[
\nabla\cdot u = 0 
\]
\[
\frac{\partial}{\partial t}(u)+\nabla\cdot (u \otimes u)- \nabla\cdot(\nu\nabla u)=-\nabla p
\]
Ahora para poderlo usar vamos a copiar el un tutorial por lo que en la terminal escribiremos lo siguiente
\begin{verbatim}
    cd $FOAM_RUN
    cp -r $FOAM_TUTORIALS/incompressible/icoFoam/cavity/cavity .
    cd cavity
\end{verbatim}

Ahora veamos que tenemos dentro de la carpeta
Primero podemos abrir la carpeta system en la cual encotraremos el archivo $\texttt{blockMeshDict}$ en el cual podemos colocar el mallado. Es importante que en la parte de $\texttt{boundary}$ tengamos claros los nombres ya que los seguiremos usando en el solver

Vamos a a parte del solver dentro de la carpeta 0 tenemos dos archivos uno $\texttt{U}$ y $\texttt{p}$, como podemos ver dentro del archivo podemos ver lo siguiente
\[
\begin{verbatim}
boudaryField
{
  movingWall
  {
    type    zeroGradient;
  }

  frontAndBack
  {
    type    empty;
  }
}
\end{verbatim}

\]

Aqui es donde pondremos los nombres que escogimos en $\texttt{boundary}$ y escogeremos alguno de los siguientes tipos para las agrupaciones de las caras.
\begin{itemize}
  \item zeroGradient:
  \item symmetryPlane:
  \item empty:
  \item fixedValue:
\end{itemize}

Ahora notemos que en la parte de hasta arriba podemos encontrar lo siguiente

$\texttt{dimensions      [0 2 -2 0 0 0 0];}$

aqui definimos las dimensiones del campo en este caso la presion kinematica $m^{2}s^{-2}$.
Debajo de este podemos ver

$\texttt{internalField uniform 0;}$

Este quiere decir que el los datos del campo interno puede ser uniforme y descritos por un solo valor o podemos usar $\texttt{nonuniform}$ donde todos los valores del campo deben de ser especificados
